\chapter{Forth}\label{cha:forth}

The machine's third language is native: no Tube, no co-processor, just
45GS10 machine code loaded at \texttt{\$4000}. It is Tali Forth 2 ---
a public-domain, ANS-flavoured Forth written for the 65C02, which this
CPU speaks as a strict superset.

\begin{type}
FORTH
\end{type}

\begin{type}
2 3 + . 5  ok
: cube dup dup * * ;  ok
7 cube . 343  ok
\end{type}

\screen{shots/forth.png}{A Forth session: the banner, arithmetic, a new
word, and the kernel disassembling itself.}

\section{The machine at your fingertips}
Forth's oldest habit is poking hardware, and this machine is one large,
friendly memory map. \verb|C@| and \verb|C!| reach every register of
Chapter~\ref{cha:io} directly:

\begin{type}
hex
D000 C@ .            read a VICKY register
D5E0 80 swap C!      silence the sound sequencer by hand
\end{type}

\section{The assembler}
Tali brings an interactive 65C02 assembler and a disassembler, kept in
because on this machine they are the point:

\begin{type}
4000 10 disasm       disassemble the Forth kernel itself
assembler-wordlist >order
here  2 lda.#  0 sta.z  drop
\end{type}

\verb|DISASM| is strictly 65C02: at a 45GS10-only opcode it prints a polite
\verb|?| and moves on.

\section{Manners}
The dictionary has about 31.7\,KB free for your words. \verb|BYE| returns
to the shell with the screen and the session exactly as you left them.

\begin{aside}
No file words yet: Forth source lives in \pth{/LANG/FORTH} but must be
typed.
An \texttt{INCLUDED} that reads through the machine's storage device is on
the list.
\end{aside}
